\section{Conclusion}
\label{sec:conclusion}

In this work, we introduced a controlled semantic perturbation framework to evaluate how structured hallucinations influence LLM-based molecular property prediction. Our results reveal a task-dependent divergence in semantic sensitivity: physicochemical prediction (BBBP) exhibits directional sensitivity to structural topic focus, while enzymatic reasoning (BACE) is robustly fragile to semantic fabrication. Crucially, we find that the semantic orientation induced by a prompt can drive predictive shifts even in the absence of factual taxonomic correctness, suggesting that hallucinations may act as a form of semantic regularization that activates latent chemical priors. 

These findings indicate that structured hallucinations are not merely errors, but informative perturbations that can be leveraged to probe the latent reasoning dynamics of LLMs. This study provides a structured methodological foundation for future work investigating the boundary between linguistic register and semantic utility in text-conditioned molecular prediction pipelines.

\subsubsection*{Future directions.} Key extensions include: (1) investigation of the distributional compression phenomenon on additional enzymatic datasets, (2) expansion of multi-model validation across diverse LLM architectures and scales, (3) larger-scale evaluation with cross-validation, and (4) comparison with fine-tuned LLM models and additional traditional ML methods to further contextualize the practical significance of the observed effects.
